\section{Literature Review}
\label{ch:litreview}
% Folded into Chapter 1 and trimmed on 20 Sep 2026 (author's instruction);
% previously a chapter of its own at roughly twice this length. The label
% ch:litreview is kept so that existing cross-references resolve; they now
% read "Section". Standing rules carried over: no italics in the running
% text, quantum references only (classical machine learning is reviewed in
% the Preliminaries), no complexity-theory or Turing-machine material.
% Citations dropped in the trim remain in the .bib.

This section surveys the literature the thesis rests on, as one argument
rather than a catalogue: the walk literature first
(Section~\ref{sec:lr-primitive}), then quantum machine learning
(Section~\ref{sec:lr-qml}), then the small body of work combining the two
(Section~\ref{sec:lr-dtqwqml}), and finally the gap the thesis addresses
(Section~\ref{sec:lr-gap}). Classical machine learning is treated in
Chapter~\ref{ch:prelim} and is not reviewed here.

\subsection{Quantum Walks as an Algorithmic Primitive}
\label{sec:lr-primitive}

The quantum walk was introduced in two forms. Aharonov, Davidovich and
Zagury~\cite{Aharonov_1993} formulated the discrete-time walk as a coin
operator on an internal degree of freedom followed by a shift conditioned on
it; Farhi and Gutmann~\cite{Farhi_1997} formulated the continuous-time walk as
evolution under the adjacency matrix of a graph, with no coin at all. The two
were later shown equivalent as models of computation~\cite{Childs2010,
strauch2006connecting}. What distinguishes either from the classical walk was
established early: on the line the walker spreads ballistically with a
bimodal limiting distribution~\cite{Ambainis_2001, Konno_2002, Grimmett_2003},
an interference effect that needs no quantisation of the walker
itself~\cite{Knight_2003}; on graphs, Aharonov et al.~\cite{Aharonov_2000}
gave the general framework and mixing-time bounds, and Childs et
al.~\cite{Childs_2001, Childs_2002} exhibited the glued-trees graph on which
the continuous-time walk beats every classical algorithm exponentially, still
the strongest separation known for a walk. Reviews are
available~\cite{Elias_2012, Kadian_2021, Xiaogang_2024}.

The coin is the feature of the discrete-time walk this thesis is concerned
with, and its literature has a consistent shape. The coin and the initial
state jointly control the distribution~\cite{Tregenna_2003}; beyond the fixed
real coins that are standard --- Hadamard, Grover and
Fourier~\cite{Komatsu_2019, Saito_2018} --- coins have been made dependent on
time~\cite{C_2006}, on the step index~\cite{Panahiyan_2018}, on
position~\cite{Ahmad_2020} and on edge weights~\cite{Wong_2017}, general
$U(2)$ coins have been parametrised and swept~\cite{Li_2012, Souza_2013,
Naves_2023}, and coin sequences have been optimised for entanglement
generation~\cite{Gratsea_2020}. This work establishes that the coin materially
changes the walk; it does not connect that freedom to a learning objective.
Alternative formulations recur later: Szegedy's quantisation of a Markov
chain~\cite{Szegedy}, equivalent to the coined walk under mild
conditions~\cite{G_2017}, the coinless staggered model~\cite{Portugal_2016},
and the lackadaisical walk with a weighted self-loop at each
vertex~\cite{G_2018}.

Search is the application in which walks are best understood. Shenvi, Kempe
and Whaley~\cite{Shenvi_2002} matched Grover's $O(\sqrt{N})$ on the hypercube
in discrete time; Childs and Goldstone~\cite{Childs_2003, Childs_2004}
established the dimensional thresholds for spatial search in continuous time;
Ambainis, Kempe and Rivosh~\cite{Ambainis_2004} and Tulsi~\cite{Tulsi_2008}
improved the two-dimensional case; Magniez et al.~\cite{Magniez_2007} gave the
general framework; and the quadratic speed-up for finding, rather than merely
detecting, a marked vertex was established in generality~\cite{Ambainis_2019,
Apers_2021}. The lackadaisical sub-literature is worth singling out because
it is the clearest demonstration that a single scalar parameter of the walk
operator decides whether a speed-up is obtained at all: only specific
self-loop weights work~\cite{Giri_2020}, the optimum shifts with the number
of marked vertices~\cite{Nahimovs_2018}, need not be uniform across
vertices~\cite{Rhodes_2019}, and may break every symmetry of the graph
without loss~\cite{Rapoza_2021}. The same pattern --- performance decided by
a few instance-dependent parameters internal to the coin --- recurs in the
learning applications below. A constant-time unstructured search using a
many-body continuous-time walk has been proposed~\cite{DalFavero_2024}; it is
a nonlinear, continuous-time result with no discrete-time analogue, and
Sections~\ref{sec:search} and~\ref{sec:res-search} record exactly what it
does and does not give.

Beyond search, Ambainis~\cite{Ambainis_2003} solved element distinctness by a
walk on the Johnson graph, matching the lower bound, and the walk has since
been applied across algebraic problems~\cite{Childs_2008}, the connection
Section~\ref{sec:shor} follows for Shor's algorithm. In optimisation, the
mixer of the quantum approximate optimisation algorithm~\cite{Farhi_2014,
Farhi_2016} is a continuous-time walk on the hypercube, which is what makes
QAOA a walk algorithm rather than one resembling it; explicitly walk-based
optimisation has been developed as the quantum walk optimisation
algorithm~\cite{Marsh_2020}, benchmarked on weighted maximum
cut~\cite{Bennett_2025} and formulated in discrete
time~\cite{Liliopoulos_2024}. Quantum random access memory has been
formulated as a walk on a binary tree~\cite{Asaka_2021, Asaka_2023}, which
bears on the data-loading question of Section~\ref{sec:lr-qml}. The
higher-level applications kept in view are entanglement distribution across
quantum networks~\cite{Chen_2024c}, multi-node teleportation~\cite{Ikken_2025},
vehicle routing under a qubit budget~\cite{Gautam_2024} and certified random
number generation~\cite{Haider_2023}: they are why the walk is treated here as
a general computational tool rather than a graph-search heuristic. Circuit
constructions for coined walks~\cite{Douglas_2007, Han_2017, Shakeel_2020,
Sato_2024} and the hardware constraints of the noisy intermediate-scale
era~\cite{Preskill_2018} bound what any of this costs to run.

\subsection{Quantum Machine Learning}
\label{sec:lr-qml}

\subsubsection{Data Encoding and Feature Maps}
Surveys of the field are available~\cite{Schuld_2015, Biamonte_2017,
Ciliberto_2018, Zeguendry_2023}, with a critical assessment by Cerezo et
al.~\cite{Cerezo_2023} and a course-level treatment by IBM
Quantum~\cite{IBMQML2026}. Early enthusiasm rested on the linear-systems
algorithm of Harrow et al.~\cite{Harrow_2008}, whose qualifications --- a
sparse well-conditioned matrix, input in superposition, and only a summary
statistic of the output --- proved restrictive: a subroutine polylogarithmic
in the input confers nothing if loading the input is linear in it. The field
has accordingly moved to algorithms that consume classical data through a
narrow interface, the feature map. Encoding a classical vector into a quantum
state fixes the function class the model can represent: Schuld, Sweke and
Meyer~\cite{Schuld_2021a} showed that a variational model with repeated
encoding is a truncated Fourier series whose spectrum is set by the encoding,
not the trainable parameters. This is why the encoding is the first design
decision in Chapter~\ref{ch:methodology}, and why a walk is a candidate
feature map at all: the walk operator is an encoding with graph structure
built into it.

\subsubsection{Quantum Neural Networks and Variational Circuits}
\label{sec:lr-qnn}
Parametrised circuits optimised by a classical routine against a measured
cost --- quantum circuit learning~\cite{Mitarai_2018}, framed as a model
class by Benedetti et al.~\cite{Benedetti_2019} and reviewed by Cerezo et
al.~\cite{Cerezo_2021} --- are called quantum neural networks or variational
circuits almost interchangeably, and are one family here. Their capacity has
been argued through the Fisher information spectrum~\cite{Abbas_2021}, and
architectures include data re-uploading~\cite{PrezSalinas_2019}, equivariant
networks~\cite{Nguyen_2024} and generative models~\cite{Gao_2018}. Two
results constrain any new proposal. McClean et al.~\cite{McClean_2018} showed
that for broad classes of random circuits the gradient vanishes exponentially
in the qubit count --- the barren plateau --- which ties depth, ansatz and
cost locality together; and Sim, Johnson and Aspuru-Guzik~\cite{Sim_2019}
showed that expressibility and entangling capability trade off against
trainability. A walk step is a particular ansatz with an entangling pattern
fixed by the graph, and it has to be judged against these measures rather
than assumed superior because the walk has good properties elsewhere.

\subsubsection{Quantum Kernel Methods}
\label{sec:lr-qkernel}
Schuld and Killoran~\cite{Schuld_2018} observed that a quantum encoding is a
feature map, so the inner product of two encoded states is a kernel and a
quantum computer can serve as a kernel evaluator; Havlicek et
al.~\cite{Havlicek_2019} implemented both a variational classifier and a
kernel estimator on hardware in one work, which makes that paper the common
ancestor of both families, and Schuld~\cite{Schuld_2021} later argued that
supervised variational models are themselves kernel methods. Overlaps are
estimated by the swap test~\cite{Buhrman_2001} or the inversion circuit, and
covariant kernels for data with group structure followed~\cite{Glick_2021}.
This line has produced the field's clearest results in both directions. Liu
et al.~\cite{Liu_2020} constructed a discrete-logarithm problem on which a
quantum kernel machine has a rigorous speed-up with classical data access.
Against that, Huang et al.~\cite{Huang_2021} showed that training data narrow
the class of problems on which an advantage can persist, and Thanasilp et
al.~\cite{Thanasilp_2024} showed that kernel values can concentrate
exponentially in the qubit count, so that distinguishing them needs
exponentially many shots. Exponential concentration is a direct design
constraint on a walk-based kernel, since a deep walk is exactly the kind of
expressive encoding that induces it.

\subsubsection{Benchmarking and the Question of Quantum Advantage}
\label{sec:lr-bench}
Schuld and Killoran~\cite{Schuld_2022} argued that quantum advantage is the
wrong organising goal for the field, the comparison being almost always
against an under-tuned baseline on a dataset chosen after the fact, and
Bowles, Ahmed and Schuld~\cite{Bowles_2024} made the point quantitatively:
across twelve quantum models on a curated benchmark collection, none beat a
well-tuned classical baseline, and apparent advantages traced to
hyperparameter and dataset selection. Their collection is the benchmark suite
this thesis uses, and their requirements --- a tuned baseline, a split fixed
before model selection, and results across seeds --- are adopted in
Chapter~\ref{ch:results}. The consequence is stated plainly: the simulations
reported later establish relative behaviour between walk configurations on
classical simulators, and no claim of quantum advantage is made anywhere in
this work.

\subsection{Quantum Walks in Quantum Machine Learning}
\label{sec:lr-dtqwqml}

The literature combining the two subjects is much smaller than either.
Schuld, Sinayskiy and Petruccione~\cite{Schuld_2014} gave the earliest
proposal, representing the firing pattern of a quantum neural network as a
walk on a graph. Dernbach et al.~\cite{Dernbach_2018, Dernbach_2019}
introduced quantum walk neural networks for graph-structured data and, in the
later work, made the coin depend on learned node features --- the closest
antecedent of this thesis. Zhu et al.~\cite{Zhu_2026} showed that a quantum
graph neural network can be realised on a single qubit by a walk, and
critical reviews of quantum graph neural networks~\cite{Ceschini_2026,
innan2026embeddingsshapegraphneural} conclude that the empirical case is not
yet made. Graph kernels built from walks predate the variational literature:
Rossi, Torsello and Hancock~\cite{Rossi_2013} built a continuous-time walk
kernel for unattributed graphs and Zhang et al.~\cite{Zhang_2022} an
R-convolution kernel from a fast discrete-time walk. These are kernels
between graphs computed from walk statistics and should be distinguished from
the quantum kernels above, which are overlaps of quantum states; the
combination --- a walk as the state-preparation circuit of a quantum kernel
--- is the construction Chapter~\ref{ch:methodology} develops, and the
literature on it is thin. Node embedding by quantum walk has attracted the
most attention of the three: QWalkVec~\cite{10.1007/978-981-97-2242-6_8}
substitutes a coined walk for the classical walk of standard embedding
methods, RED~\cite{Wang_2021} learns role embeddings by discrete-time walk,
and GraphQWalk~\cite{Liu_2026} the structural analogue in continuous time;
walks have also been applied to community detection~\cite{Mukai_2020,
He_2025} and to centrality~\cite{Wang_2025, Berry_2010}. The reported gains
in this sub-literature are typically small, obtained on a handful of small
graphs, and sensitive to the truncation length and the downstream
classifier. Node embedding is recorded here as related work only and is not
pursued in this thesis; the same applies to walks with
memory~\cite{Flitney_2003, McGettrick_2009, Shakeel_2019,
PhysRevA.106.062215}, proposed as the quantum counterpart of the
second-order classical walk.

\subsection{The Gap Addressed by This Thesis}
\label{sec:lr-gap}

Three observations follow from the survey and together define the gap. First,
the walk literature establishes repeatedly that the performance of a
walk-based algorithm is governed by a few parameters internal to the coin
whose correct values are instance-dependent; the lackadaisical self-loop
weight is the clearest case~\cite{Nahimovs_2018, Rhodes_2019, Rapoza_2021}.
The coin is known there to be a design space, not a convention. Second, the
walk-based learning literature does not treat it as one: with the partial
exception of Dernbach et al.~\cite{Dernbach_2019}, walk-based learning models
fix the coin to Hadamard, Grover or Fourier and report results for that
choice, although position-dependent~\cite{Ahmad_2020},
step-dependent~\cite{Panahiyan_2018} and general $U(2)$~\cite{Li_2012,
Souza_2013} coins are all standard in the walk literature, and none has been
evaluated against a learning objective. Third, the benchmarking literature
shows why this has gone unnoticed: where a walk-based model beats a weak
baseline there is no reason to ask whether another coin would have done
better, and against tuned baselines walk-based models have not so far
won~\cite{Bowles_2024}. The coin has been neither optimised nor ruled out.

The thesis accordingly asks a narrow question. Taking the discrete-time coined
walk as a computational primitive and instantiating it in the mainstream
families above --- quantum neural networks with variational classifiers, and
quantum kernel methods --- what does treating the coin as a design space buy?
Two axes are varied: whether the coin depends on local graph structure and on
the step index, and whether its entries may be complex. The question is
empirical, answered on simulators against tuned classical baselines under the
reporting discipline of Bowles et al.~\cite{Bowles_2024}, and a negative
answer on either axis is a result and is reported as one.

